\section{Approcci multi-canale e tecniche di selezione del link}
\label{sec:sota_multichannel}

La complementarità tra tecnologie (ad es.\ radio dedicata, cellulare e VLC) ha portato a un crescente interesse verso architetture multi-canale, in cui l’affidabilità complessiva si ottiene tramite \emph{diversity}, ridondanza o selezione dinamica del link in funzione delle condizioni del canale e delle esigenze del controllo \cite{hardes2019communication,ishihara2015improving}. Tuttavia, l’uso simultaneo o alternato di più interfacce non è banale: scelte ingenue di switching possono introdurre instabilità o violazioni dei vincoli di sicurezza, perché la logica di controllo deve essere coerente con l’interfaccia attiva e con la qualità informativa del flusso di messaggi \cite{giordano2019joint,segata2023multi-technology}. La letteratura più recente affronta quindi il problema con protocolli e procedure di \emph{fallback/recovery} che coordinano, in modo distribuito, transizioni sicure tra controllori e tra tecnologie, fino a modalità non cooperative quando necessario \cite{segata2023multi-technology}. Questa sezione inquadra le principali strategie multi-canale e le metriche comunemente utilizzate (ad es.\ PDR e indicatori di qualità del link), evidenziando i nodi progettuali che la parte sperimentale della tesi intende analizzare e validare.